% Caption for Fig Threads (fig-threads.tex), thread-scaling appendix. Numbers from the
% threads-2885ee8-da1041e8 sweep.
\caption{Thread scaling of matching statistics (MS) queries on the HPRC2 index (Human472, 472
haplotypes). Left: query speed (Mbases/s) against the number of threads for the ms2 modes min, psi,
phi and dual and for ropebwt3 \texttt{-l1}, which reports every MS value. Right: speedup over the
same configuration at one thread; the dashed line marks linear scaling up to 32 threads. Query
sets: Short reads (Illumina, ACGT-only), 7,982,098 reads, 997.8 Mbases; Nanopore (HG002 Q28),
43,346 reads, 1,000.0 Mbases. Each set extends the corresponding 60 Mbase set of
Fig.~\ref{fig:ms}, which is its prefix. All runs used one AWS i4i.metal instance (two sockets of
32 cores each). With $T \le 32$ threads, each thread was bound to its own physical core of the
first socket. At $T = 64$ (gray band), two threads shared each of those 32 cores through
hyperthreading, also called simultaneous multithreading (SMT). Memory was bound to the same
socket in every run. ms2 split the reads into 2 MiB blocks and ropebwt3 ran with
\texttt{-t}~$T$. Speed is query time only; index load time is excluded. Each value is the median
of three runs, with the page cache dropped before each run. The output of every run matched the
ropebwt3 truth at every position, and each configuration's output was byte-identical at all
seven thread counts. ms2 commit 2885ee8; ropebwt3 commit b408339.}
\label{fig:threads}
